% ---------------------------------------------------------------------------
%  Preámbulo LaTeX del manual (formato PDF)
%  Se inserta en la cabecera del documento generado por Quarto.
% ---------------------------------------------------------------------------

% ----------------------------- Colores de marca ----------------------------
\definecolor{eaNaranja}{HTML}{F39200}
\definecolor{eaNaranjaClaro}{HTML}{FFD79A}
\definecolor{eaVerde}{HTML}{4F8A2B}
\definecolor{eaVerdeOscuro}{HTML}{27451C}
\definecolor{eaGris}{HTML}{4A5568}
\definecolor{eaGrisClaro}{HTML}{F2F4EE}
\definecolor{eaRojo}{HTML}{C0392B}
\definecolor{eaAmbar}{HTML}{E6A700}

% ----------------------- Dibujo de la portada (TikZ) -----------------------
\usepackage{tikz}
\usetikzlibrary{calc}

% ------------------------------- Tipografía --------------------------------
\usepackage[default,scale=0.96]{sourcesanspro}
\usepackage[T1]{fontenc}
\usepackage{microtype}

% ------------------------- Encabezados y pies de página --------------------
\usepackage[automark,headsepline=0.6pt,footsepline=0.6pt]{scrlayer-scrpage}
\pagestyle{scrheadings}
\clearpairofpagestyles
\automark[section]{chapter}
\ihead{\footnotesize\color{eaGris}\headmark}
\ohead{}
\ofoot*{\footnotesize\color{eaVerdeOscuro}\pagemark}
\ifoot*{\footnotesize\color{eaGris}Manual de uso · Mapa de Conflictos Ambientales}
\setkomafont{pageheadfoot}{\normalfont\sffamily}
\setkomafont{headsepline}{\color{eaNaranja}}
\setkomafont{footsepline}{\color{eaNaranja}}

% ---------------------- Estilo de capítulos y secciones --------------------
\addtokomafont{disposition}{\sffamily}
\setkomafont{chapter}{\color{eaVerdeOscuro}\sffamily\bfseries\huge}
\setkomafont{section}{\color{eaVerde}\sffamily\bfseries\Large}
\setkomafont{subsection}{\color{eaVerdeOscuro}\sffamily\bfseries\large}
\setkomafont{subsubsection}{\color{eaGris}\sffamily\bfseries\normalsize}
\setkomafont{part}{\color{eaVerdeOscuro}\sffamily\bfseries\Huge}
\setkomafont{partnumber}{\color{eaNaranja}\sffamily\bfseries\Huge}

% Número de capítulo grande, en naranja, con filete bajo el título
\renewcommand*{\chapterformat}{%
  \mbox{\color{eaNaranja}\fontsize{54}{54}\selectfont\bfseries\thechapter\enskip}%
}
\renewcommand*{\chapterlinesformat}[3]{%
  \parbox{\textwidth}{\raggedright #2#3%
    \par\nobreak\vspace{-0.2\baselineskip}%
    {\color{eaNaranja}\rule{\linewidth}{1.2pt}}}%
}

% ------------------------- Índice general (TOC) ----------------------------
\usepackage{tocbasic}
\DeclareTOCStyleEntry[entryformat=\bfseries\color{eaVerdeOscuro}]{tocline}{chapter}

% ------------------------------- Enlaces -----------------------------------
% (los colores de los enlaces se fijan en _quarto.yml mediante colorlinks)

% --------------------------- Figuras y capturas ----------------------------
\usepackage{float}
\usepackage[export]{adjustbox}       % permite bordes en \includegraphics
\makeatletter
% Marco gris claro para las capturas de pantalla
\newcommand{\captura}[2][\textwidth]{%
  \begin{center}
    \fbox{\includegraphics[width=\dimexpr#1-2\fboxsep-2\fboxrule\relax]{#2}}
  \end{center}}
\makeatother

% ---------------------- Elementos propios del manual -----------------------
% Equivalentes en PDF de las clases CSS usadas en la versión web.
\usepackage{tcolorbox}
\tcbuselibrary{skins,breakable}

% Teclas e interfaz:  [Guardar]{.ui}
\newcommand{\ui}[1]{%
  \tcbox[on line, boxsep=1pt, left=2pt, right=2pt, top=1pt, bottom=1pt,
         colframe=eaVerdeOscuro!55, colback=white, boxrule=0.5pt, arc=2pt]%
        {\sffamily\small\bfseries\color{eaVerdeOscuro}#1}}

% Rutas de menú:  [Inicio > Archivo]{.ruta}
\newcommand{\ruta}[1]{{\sffamily\bfseries\color{eaNaranja}#1}}

% Etiquetas de gravedad
\newcommand{\etiqueta}[2]{%
  \tcbox[on line, boxsep=1pt, left=3pt, right=3pt, top=0pt, bottom=0pt,
         colframe=#1, colback=#1, arc=4pt]{\sffamily\scriptsize\bfseries\color{white}\MakeUppercase{#2}}}
\newcommand{\gravedadalta}[1]{\etiqueta{eaRojo}{#1}}
\newcommand{\gravedadmedia}[1]{\etiqueta{eaAmbar}{#1}}
\newcommand{\gravedadbaja}[1]{\etiqueta{eaVerde}{#1}}

% Las divisiones con clase (::: {.pasos}) se convierten en entornos homónimos
\newenvironment{pasos}{\medskip}{\medskip}
\newenvironment{tarjeta}
  {\begin{tcolorbox}[breakable, colback=eaGrisClaro, colframe=eaNaranja,
                     boxrule=0pt, toprule=2.5pt, arc=2pt, left=8pt, right=8pt]}
  {\end{tcolorbox}}

% --------------------- Caracteres especiales (Unicode) ---------------------
\usepackage{amssymb}
\DeclareUnicodeCharacter{2713}{\ensuremath{\checkmark}}   % marca de verificación
\DeclareUnicodeCharacter{2717}{\ensuremath{\times}}
\DeclareUnicodeCharacter{203A}{\guilsinglright}           % flecha de ruta de menú
\DeclareUnicodeCharacter{2039}{\guilsinglleft}
\DeclareUnicodeCharacter{00B7}{\textperiodcentered}
\DeclareUnicodeCharacter{2192}{\ensuremath{\rightarrow}}
\DeclareUnicodeCharacter{2026}{\textellipsis}
\DeclareUnicodeCharacter{2212}{\textminus}
\DeclareUnicodeCharacter{2715}{\ensuremath{\times}}

% ------------------------------ Varios -------------------------------------
\usepackage{enumitem}
\setlist{itemsep=2pt, parsep=0pt, topsep=4pt}
\renewcommand{\arraystretch}{1.15}
